\Qhead{6.1}{Задачи оптимального приёма сигнала. Статистическая проверка гипотез}{Задачи оптимального приема сигнала. Апостериорная плотность вероятности. Функция правдоподобия. Статистическая проверка гипотез. Критерии Байеса, Неймана-Пирсона и Вальда проверки гипотез.}
\Subfile{6.1.1}{sections/06_detection/q01/s01}
\Subfile{6.1.2}{sections/06_detection/q01/s02}

\recap{Сначала задают модель, априорные сведения и цену ошибки. Для бинарного решения центрально отношение правдоподобий; разные критерии выбирают порог и момент остановки по разным требованиям.}
\board{Байес; гауссово отношение; три критерия и две границы Вальда.}
\oral{Почему максимум правдоподобия и максимум апостериорной вероятности могут дать разные ответы?}
\bait{Можно оставить случайность времени Вальда: экономия достигается тем, что убедительные данные завершают наблюдение раньше.}
\sources{Университетский курс математической статистики~\cite{msustats}, правдоподобие и Нейман–Пирсон; MIT 6.432~\cite{wornell}, бинарное обнаружение. Формулы байесовского риска выведены здесь с явно заданными индексами цен; границы Вальда — приближённые.}
